\ResearchStyle
\PartDivider{PART II}{Research Problem and Gap}{Reconstruct the exact unfairness problem that motivates CATRA.}

\begin{frame}{Paper Context and Scope}
\PaperClaim{The paper studies contention between stations and between flows in IEEE 802.11 multi-hop ad hoc networks.}
\begin{columns}[T,onlytextwidth]
\column{0.48\textwidth}
\textbf{Included}
\begin{itemize}
\item Shared IEEE 802.11 medium.
\item Multi-hop forwarding.
\item TCP flow contention.
\item MAC and transport feedback.
\end{itemize}
\column{0.48\textwidth}
\textbf{Explicitly excluded}
\begin{itemize}
\item Mobility-induced unfairness.
\item Multi-rate operation.
\item Multi-channel operation.
\end{itemize}
\end{columns}
\end{frame}

\begin{frame}{Problem 1: MAC-Layer Unfairness}
\begin{center}
\resizebox{0.96\textwidth}{!}{%
\begin{tikzpicture}[node distance=0.62cm]
\node[draw=ResearchColor,rounded corners,fill=CardBg] (beb) {BEB chooses CW};
\node[draw=ResearchColor,rounded corners,fill=CardBg,right=of beb] (state) {Uses contention state};
\node[draw=WarningColor,rounded corners,fill=CardBg,right=of state] (blind) {No flow awareness};
\node[draw=WarningColor,rounded corners,fill=CardBg,right=of blind] (unfair) {Unequal opportunities};
\draw[-{Latex[length=2mm]}] (beb)--(state);
\draw[-{Latex[length=2mm]}] (state)--(blind);
\draw[-{Latex[length=2mm]}] (blind)--(unfair);
\end{tikzpicture}
}
\end{center}
\vspace{0.5cm}
\begin{itemize}
\item BEB appears symmetric but multi-hop contention domains are asymmetric.
\item A forwarding station may compete with different sets of neighbors on each side.
\item The standard backoff rule does not explicitly represent how many flows a station must carry.
\item A disadvantaged station can repeatedly lose access even when it forwards critical traffic.
\end{itemize}
\end{frame}

\begin{frame}{Problem 2: TCP-Layer Unfairness}
\begin{center}
\resizebox{0.96\textwidth}{!}{%
\begin{tikzpicture}[node distance=0.58cm]
\node[draw=ResearchColor,rounded corners,fill=CardBg] (grow) {$cwnd$ grows};
\node[draw=ResearchColor,rounded corners,fill=CardBg,right=of grow] (packets) {More packets};
\node[draw=WarningColor,rounded corners,fill=CardBg,right=of packets] (relay) {Relay queues};
\node[draw=WarningColor,rounded corners,fill=CardBg,right=of relay] (cont) {More contention};
\node[draw=WarningColor,rounded corners,fill=CardBg,right=of cont] (fair) {Fairness loss};
\draw[-{Latex[length=2mm]}] (grow)--(packets);
\draw[-{Latex[length=2mm]}] (packets)--(relay);
\draw[-{Latex[length=2mm]}] (relay)--(cont);
\draw[-{Latex[length=2mm]}] (cont)--(fair);
\end{tikzpicture}
}
\end{center}
\vspace{0.45cm}
\begin{itemize}
\item Conventional TCP controls an end-to-end path but does not directly observe local wireless contention.
\item A large window creates many competing packets at intermediate relays.
\item Different paths and start times can produce unequal ACK clocks and unequal window growth.
\item Transport behavior that is valid on wired paths can overload a shared multi-hop medium.
\end{itemize}
\end{frame}

\begin{frame}{Problem Statement}

\begin{columns}[T,onlytextwidth]
\column{0.48\textwidth}
\textbf{MAC-layer unfairness}
\begin{itemize}
    \item BEB adjusts CW mainly from contention/collision state.
    \item It does not explicitly consider neighboring-flow demand.
    \item Some stations can gain more channel access opportunities.
\end{itemize}

\column{0.48\textwidth}
\textbf{TCP-layer unfairness}
\begin{itemize}
    \item TCP cwnd can grow too aggressively.
    \item More packets contend for the same medium.
    \item Throughput and per-flow fairness degrade.
\end{itemize}
\end{columns}

\vspace{0.45cm}

\begin{center}
\begin{tikzpicture}[node distance=0.8cm]

    \node[
        draw=ResearchColor,
        rounded corners,
        fill=CardBg,
        minimum width=2.4cm,
        minimum height=0.75cm
    ] (mac) {MAC contention};

    \node[
        draw=ResearchColor,
        rounded corners,
        fill=CardBg,
        minimum width=2.4cm,
        minimum height=0.75cm,
        right=1.1cm of mac
    ] (tcp) {TCP sending rate};

    \node[
        draw=WarningColor,
        rounded corners,
        fill=CardBg,
        minimum width=3.3cm,
        minimum height=0.75cm,
        right=1.1cm of tcp
    ] (fair) {Fairness degradation};

    \draw[
        -{Latex[length=2.5mm]},
        draw=ResearchColor,
        line width=0.8pt
    ]
    (mac) -- (tcp);

    \draw[
        -{Latex[length=2.5mm]},
        draw=WarningColor,
        line width=0.8pt
    ]
    (tcp) -- (fair);

\end{tikzpicture}
\end{center}
\end{frame}

\begin{frame}[t]{Research Gap}
% =========================================================
% Key point
% =========================================================
\begin{center}
\begin{tikzpicture}

\node[
    draw=ResearchColor,
    rounded corners=4pt,
    fill=CardBg,
    text width=13.2cm,
    align=center,
    inner xsep=8pt,
    inner ysep=6pt
] {
    \large\bfseries
    Existing solutions improve either MAC access or TCP sending
    behavior, but do not sufficiently address station-level and
    flow-level contention together.
};

\end{tikzpicture}
\end{center}

\vspace{0.10cm}

% =========================================================
% Research gap diagram
% =========================================================
\begin{center}

\begin{tikzpicture}[
    node distance=0.55cm and 0.65cm,
    approach/.style={
        draw=ResearchColor,
        rounded corners=4pt,
        fill=CardBg,
        text width=2.8cm,
        align=center,
        minimum height=0.82cm,
        inner sep=4pt,
        font=\small
    },
    gap/.style={
        draw=MethodColor,
        line width=1pt,
        rounded corners=4pt,
        fill=CardBg,
        text width=4.2cm,
        align=center,
        minimum height=0.90cm,
        inner sep=5pt,
        font=\small
    }
]

% ---------------------------------------------------------
% Existing approaches
% ---------------------------------------------------------

\node[approach] (mac) {
    MAC-only\\
    approaches
};

\node[
    approach,
    right=0.65cm of mac
] (tcp) {
    TCP-only\\
    approaches
};

\node[
    approach,
    right=0.65cm of tcp
] (cross) {
    Previous cross-layer\\
    approaches
};

% ---------------------------------------------------------
% Join point
% ---------------------------------------------------------

\coordinate[
    below=0.38cm of tcp
] (join);

% ---------------------------------------------------------
% Research gap / need
% ---------------------------------------------------------

\node[
    gap,
    below=0.42cm of join
] (catra) {
    \textbf{Research need}\\[-1pt]
    Joint MAC + TCP fairness control
};

% ---------------------------------------------------------
% Convergence lines
% ---------------------------------------------------------

\draw[
    draw=ResearchColor,
    line width=0.75pt
]
(mac.south) |- (join);

\draw[
    draw=ResearchColor,
    line width=0.75pt
]
(tcp.south) -- (join);

\draw[
    draw=ResearchColor,
    line width=0.75pt
]
(cross.south) |- (join);

% ---------------------------------------------------------
% Final arrow
% ---------------------------------------------------------

\draw[
    -{Latex[length=2.3mm]},
    draw=MethodColor,
    line width=0.9pt
]
(join) -- (catra.north);

\end{tikzpicture}

\end{center}

\end{frame}

\begin{frame}{Research Need: Two-Stage Control}
\begin{center}
\begin{tikzpicture}[node distance=0.8cm]
\node[draw=ResearchColor,rounded corners,fill=CardBg,text width=4.2cm,align=center,minimum height=1.1cm] (s1)
{Stage 1\\Determine a suitable CW for each station};
\node[draw=ResearchColor,rounded corners,fill=CardBg,text width=4.2cm,align=center,minimum height=1.1cm,right=of s1] (s2)
{Stage 2\\Regulate the sending rate of each TCP flow};
\draw[-{Latex[length=2.8mm]},draw=ResearchColor,line width=1.2pt] (s1)--(s2);
\end{tikzpicture}
\end{center}
\vspace{0.5cm}
\begin{itemize}
\item Stage 1 targets fair channel access between stations.
\item Stage 2 targets fair bandwidth sharing between end-to-end flows.
\item Both stages use channel-utilization information collected at MAC/PHY.
\end{itemize}
\end{frame}

\begin{frame}{Research Questions}

\begin{enumerate}
    \item Can adaptive CW control based on measured channel utilization improve per-station fairness in multi-hop IEEE 802.11 networks?

    \item Can MAC/PHY information be used to dynamically regulate TCP sending rate and improve per-flow fairness?

    \item Can both mechanisms be combined while maintaining reasonably good channel utilization and throughput?
\end{enumerate}
\end{frame}


\begin{frame}{Main Contributions of the Paper}
\begin{enumerate}
\item Define station-level and flow-level Fair Bandwidth Ratios.
\item Estimate real channel utilization from packet airtime.
\item Adapt the IEEE 802.11 contention window using $RBR_S/FBR_S$.
\item Pass MAC/PHY utilization information to the transport layer.
\item Adapt TCP generation delay and $cwnd$ using $RBR_f/FBR_f$.
\end{enumerate}
\vspace{0.35cm}
\KeyPoint{CATRA is a feedback-control design with separate actuators for station contention and flow demand.}
\end{frame}
